\chapter{Ce que veut dire « apprendre » pour une machine}
\label{chap:apprendre}

\begin{quote}\itshape
Une image, nous venons de le voir, n'est rien d'autre qu'un grand tas de
nombres. Mais un tas de nombres ne fait rien tout seul. Dans ce chapitre,
nous allons voir comment une machine part d'un réglage pris au hasard, se
trompe, remarque qu'elle s'est trompée, et se corrige toute seule. C'est
cela, et rien d'autre, qu'on appelle « apprendre ».
\end{quote}

\begin{prerequis}
Le chapitre~\ref{chap:image-video}, et seulement lui. Il faut savoir qu'une
image est une grille de nombres, et savoir ce qu'est la forme d'un tenseur.
Pour le reste, il suffit de savoir additionner, multiplier, et diviser. Nous
n'utiliserons aucune formule que nous n'aurons pas d'abord dessinée.
\end{prerequis}

Quand on dit qu'une machine « apprend », on imagine souvent quelque chose de
mystérieux, presque magique. La réalité est beaucoup plus simple, et
beaucoup plus rassurante : une machine qui apprend, c'est une machine qui
essaie, qui se trompe, qui regarde de combien elle s'est trompée, et qui
recommence un tout petit peu mieux. Des milliers de fois. Des millions de
fois.

Tout ce chapitre raconte cette boucle-là. Nous allons la suivre pas à pas
sur un exemple minuscule, que nous calculerons entièrement à la main, avant
de vérifier sur une vraie machine que nos calculs de papier tombent juste.

\section{Une machine qui apprend, c'est une machine à boutons}

\begin{enclair}
Imaginez un appareil avec quelques boutons à tourner. Vous lui posez une
question, il donne une réponse. Si la réponse est mauvaise, vous tournez un
peu les boutons et vous redemandez. Apprendre, pour une machine, c'est
exactement cela : trouver la bonne position des boutons.
\end{enclair}

Vous avez sûrement déjà réglé la température d'une douche. Il y a un
robinet chaud et un robinet froid. Vous les tournez un peu, vous mettez la
main sous l'eau, c'est trop froid, alors vous tournez encore. Au bout de
quelques essais, l'eau est parfaite. Personne ne vous a donné la position
exacte des robinets : vous l'avez trouvée en essayant et en corrigeant.

Une machine qui apprend fait la même chose, sauf qu'elle a beaucoup plus de
robinets, et qu'elle ne trempe pas la main dans l'eau : elle calcule.

\begin{definition}[Le modèle][modele]
Un \emph{modèle} est une machine à calculer qui prend quelque chose en
entrée et rend une réponse en sortie. Ce qu'elle rend ne dépend pas
seulement de l'entrée : cela dépend aussi de la position de ses boutons. Le
même modèle, avec des boutons réglés autrement, donnera une autre réponse à
la même question.
\end{definition}

\begin{definition}[Le paramètre, ou poids][parametre]
Un \emph{paramètre} est un de ces boutons. C'est simplement un nombre, rangé
dans la mémoire de la machine, que l'on a le droit de changer. On l'appelle
aussi un \emph{poids}, parce qu'il décide de l'importance que la machine
donne à telle ou telle information. Régler un modèle, c'est choisir la
valeur de tous ses paramètres.
\end{definition}

La figure~\ref{fig:boutons} montre cette idée dans son entier : une entrée
qui arrive, des boutons qu'on peut tourner, une réponse qui sort.

\begin{schema}{Un modèle, c'est une boîte à calculer munie de boutons. On
lui donne une entrée par la gauche, elle rend une réponse par la droite.
Entre les deux, ce sont les boutons — les paramètres — qui décident du
résultat. Tourner les boutons change la réponse sans rien changer à la
question.}[boutons]
  \begin{scope}[x=1cm, y=1cm]
    % --- la boîte du modèle ---
    \draw[draw=mbBlue, fill=mbBlue!8, line width=1pt, rounded corners=4pt]
          (0,0) rectangle (5,2.5);
    \node[font=\sffamily\small\bfseries, text=mbBlue] at (2.5,0.45)
         {le modèle};
    % --- les trois boutons ---
    % Les boutons sont numérotés, et non nommés a, b, c : le modèle de la
    % page suivante n'en a que deux, et il ne faut pas que le lecteur croie
    % qu'il en manque un.
    \foreach \mbk/\mbang/\mbnom in {0/55/1, 1/-30/2, 2/120/3}{
      \begin{scope}[shift={(1.3+\mbk*1.2, 1.6)}]
        \draw[draw=mbSlate!65, fill=white, line width=0.9pt] (0,0) circle (0.34);
        \draw[draw=mbAmber!85!black, line width=1.3pt] (0,0) -- (\mbang:0.27);
        \fill[mbInk] (0,0) circle (0.045);
        \node[font=\sffamily\scriptsize, text=mbSlate] at (0,-0.62) {\mbnom};
      \end{scope}
    }
    \node[font=\sffamily\scriptsize\itshape, text=mbAmber!75!black]
         at (2.5,2.85) {les boutons, c'est-à-dire les paramètres};
    % --- entrée ---
    \draw[fleche] (-2.1,1.25) -- (-0.12,1.25);
    \node[font=\sffamily\small, text=mbInk, align=center] at (-2.9,1.25)
         {la question\\ $x$};
    % --- sortie ---
    \draw[fleche accent] (5.12,1.25) -- (7.1,1.25);
    \node[font=\sffamily\small, text=mbInk, align=center] at (7.95,1.25)
         {la réponse\\ $\pred{y}$};
  \end{scope}
\end{schema}

Prenons tout de suite un exemple, et gardons-le jusqu'à la fin du chapitre.

\begin{exemple}[La machine à bonbons]
Dans le hall d'une école, il y a une machine à bonbons. On y glisse des
pièces, et elle rend des bonbons. Trois enfants ont essayé et ont noté ce
qu'ils ont obtenu :

\begin{center}
\begin{tabular}{ccc}
\toprule
\textbf{Enfant} & \textbf{Pièces glissées} & \textbf{Bonbons reçus} \\
\midrule
Léa   & 1 & 2 \\
Sami  & 2 & 5 \\
Nina  & 3 & 8 \\
\bottomrule
\end{tabular}
\end{center}

Nous aimerions une machine à calculer qui devine le nombre de bonbons à
partir du nombre de pièces. Appelons $x$ le nombre de pièces et $\pred{y}$
le nombre de bonbons deviné. Le modèle le plus simple du monde s'écrit
\[
  \pred{y} \;=\; a \times x \;+\; b .
\]
Il a exactement deux boutons : $a$ et $b$. Le bouton $a$ dit combien de
bonbons on gagne par pièce ; le bouton $b$ ajoute ou enlève un nombre fixe
de bonbons, quel que soit le nombre de pièces.
\end{exemple}

\begin{intuition}
Le bouton $a$, c'est l'inclinaison de la droite : plus il est grand, plus la
machine est généreuse à chaque pièce. Le bouton $b$, c'est le cadeau de départ : il
décale toute la réponse vers le haut ou vers le bas, même quand on n'a mis
aucune pièce.
\end{intuition}

Vous avez peut-être déjà trouvé la règle de tête : trois bonbons par pièce,
moins un bonbon que la machine garde toujours pour elle. Autrement dit
$a = 3$ et $b = -1$. Vérifions sur les trois enfants : $3 \times 1 - 1 = 2$
pour Léa, $3 \times 2 - 1 = 5$ pour Sami, $3 \times 3 - 1 = 8$ pour Nina.
Tout tombe juste.

\begin{piege}
Nous connaissons la bonne réponse, mais \textbf{la machine, elle, ne la
connaît pas}. C'est tout l'intérêt de l'exercice : nous allons la laisser
partir de boutons réglés à zéro, complètement à côté de la plaque, et la
regarder retrouver toute seule le $3$ et le $-1$. Connaître la réponse à
l'avance nous sert seulement à vérifier qu'elle ne se trompe pas de chemin.
\end{piege}

\begin{definition}[La prédiction][prediction]
Une \emph{prédiction} est la réponse que le modèle donne avec les boutons
tels qu'ils sont réglés en ce moment. On la note $\pred{y}$, avec un petit
chapeau, pour bien la distinguer de la vraie valeur $y$ que l'on a observée
dans la réalité.
\end{definition}

\begin{exemple}[Une première prédiction, très mauvaise]
Réglons les deux boutons à zéro : $a = 0$ et $b = 0$. Pour Léa, qui a mis
une pièce :
\[
  \pred{y} \;=\; 0 \times 1 + 0 \;=\; 0 .
\]
La machine annonce zéro bonbon, alors que Léa en a reçu deux. Elle se trompe
donc de deux bonbons. Pour Sami elle annonce encore zéro, alors qu'il en a
reçu cinq. Et pour Nina, zéro contre huit. Notre modèle est pour l'instant
parfaitement nul, et c'est normal : nous n'avons rien réglé du tout.
\end{exemple}

Nous avons une machine, des boutons, et des réponses fausses. Il nous
manque quelque chose d'essentiel : un moyen de dire \emph{à quel point}
elles sont fausses.

\section{Se tromper, et mesurer de combien}

\begin{enclair}
Pour se corriger, il ne suffit pas de savoir qu'on s'est trompé. Il faut un
nombre qui dise si l'on s'est trompé beaucoup ou un tout petit peu. Ce
nombre s'appelle l'erreur.
\end{enclair}

Reprenons Léa. La machine a dit $0$, la réalité était $2$. L'écart entre les
deux vaut $0 - 2 = -2$. Nous pourrions nous arrêter là et additionner les
écarts des trois enfants. Mais un problème apparaît aussitôt.

\begin{piege}
Additionner les écarts tels quels ne marche pas. Imaginez une machine qui se
trompe de $+5$ bonbons pour Léa et de $-5$ bonbons pour Sami. En
additionnant, on obtient $5 + (-5) = 0$ : on croirait que la machine est
parfaite, alors qu'elle s'est lourdement trompée deux fois ! Les erreurs se
sont annulées entre elles.
\end{piege}

Il faut donc empêcher les écarts de se compenser. La solution retenue
presque partout est de \textbf{mettre chaque écart au carré}, c'est-à-dire de
le multiplier par lui-même. Un carré est toujours positif : $(+5)^2 = 25$ et
$(-5)^2 = 25$ aussi. Plus rien ne s'annule.

\begin{intuition}
Mettre au carré a un deuxième avantage, qui est même le principal : cela
punit très fort les grosses erreurs. Se tromper de $1$ coûte $1$. Se tromper
de $2$ ne coûte pas $2$ mais $4$. Se tromper de $10$ coûte $100$. La machine
va donc s'occuper en priorité de ses plus grosses bêtises, ce qui est
exactement ce qu'on veut.
\end{intuition}

\begin{definition}[L'erreur, ou perte][perte]
L'\emph{erreur} est un seul nombre qui résume à quel point le modèle se
trompe sur l'ensemble des exemples. On l'appelle aussi la \emph{perte}, et
on la note souvent $\mathcal{L}$, de l'anglais \emph{loss}. On la calcule en
mettant chaque écart au carré, puis en prenant la moyenne :
\[
  \mathcal{L} \;=\; \frac{1}{N}\sum_{i=1}^{N}
                    \bigl(\pred{y}_i - y_i\bigr)^{2},
\]
où $N$ est le nombre d'exemples, $\pred{y}_i$ la réponse du modèle sur
l'exemple numéro $i$, et $y_i$ la vraie valeur. Une erreur de $0$ veut dire
que le modèle tombe juste partout. Plus le nombre est grand, plus le modèle
est mauvais.
\end{definition}

\begin{notation}[Lire le signe somme][somme]
Le grand symbole $\sum$ se lit « somme de ». L'écriture
$\sum_{i=1}^{3} z_i$ veut simplement dire $z_1 + z_2 + z_3$ : on fait
prendre à $i$ toutes les valeurs de $1$ à $3$, et on additionne tout. C'est
une façon courte d'écrire une addition longue, rien de plus. Si l'addition
ne comporte que trois termes, vous avez parfaitement le droit de les écrire
à la main : c'est ce que nous ferons.
\end{notation}

\begin{exemple}[L'erreur de notre machine réglée à zéro]
Avec $a = 0$ et $b = 0$, les trois prédictions valent $0$. Calculons les
écarts, puis leurs carrés :

\begin{center}
\begin{tabular}{lccccc}
\toprule
\textbf{Enfant} & $x$ & $y$ & $\pred{y}$ & \textbf{écart } $\pred{y}-y$
& \textbf{carré} \\
\midrule
Léa  & 1 & 2 & 0 & $-2$ & $4$  \\
Sami & 2 & 5 & 0 & $-5$ & $25$ \\
Nina & 3 & 8 & 0 & $-8$ & $64$ \\
\bottomrule
\end{tabular}
\end{center}

La somme des carrés vaut $4 + 25 + 64 = 93$. Comme il y a trois exemples, la
moyenne vaut
\[
  \mathcal{L} \;=\; \frac{93}{3} \;=\; 31 .
\]
Retenez ce nombre : \textbf{l'erreur de départ vaut exactement $31$}. Nous
allons passer le reste du chapitre à la faire descendre.
\end{exemple}

La figure~\ref{fig:erreur-droite} donne à voir ce que ces écarts
représentent : ce sont des distances verticales entre les vrais points et la
droite que dessine le modèle.

\begin{schema}{Les trois enfants sont les points bleus. La droite en trait
plein est la réponse du modèle. Ici, ses deux boutons sont encore à zéro :
elle est couchée sur l'axe et annonce zéro bonbon à tout le monde. Les
traits rouges verticaux sont les écarts, c'est-à-dire de combien la machine
s'est trompée sur chaque enfant : $2$, $5$ et $8$ bonbons. Régler les
boutons, ce sera faire pivoter et monter cette droite jusqu'à ce que les
traits rouges soient les plus courts possible.}[erreur-droite]
  \begin{scope}[x=1.25cm, y=0.36cm]
    % --- axes ---
    \draw[->, draw=mbSlate, line width=0.7pt] (0,0) -- (3.9,0)
          node[right, font=\sffamily\scriptsize, text=mbSlate] {pièces $x$};
    \draw[->, draw=mbSlate, line width=0.7pt] (0,0) -- (0,10.2)
          node[above, font=\sffamily\scriptsize, text=mbSlate] {bonbons};
    \foreach \mbk in {1,2,3}{
      \draw[draw=mbSlate!60, line width=0.6pt] (\mbk,-0.35) -- (\mbk,0.35);
      \node[font=\sffamily\scriptsize, text=mbSlate] at (\mbk,-1.1) {\mbk};
    }
    \foreach \mbk in {2,4,6,8}{
      \draw[draw=mbSlate!60, line width=0.6pt] (-0.07,\mbk) -- (0.07,\mbk);
      \node[font=\sffamily\scriptsize, text=mbSlate, left] at (-0.09,\mbk) {\mbk};
    }
    % --- la droite du modèle, boutons à zéro : y = 0 ---
    \draw[draw=mbAmber!85!black, line width=1.4pt] (0,0) -- (3.6,0);
    \draw[draw=mbAmber!65, line width=0.5pt] (0.45,-1.55) -- (0.45,-0.15);
    \node[font=\sffamily\scriptsize, text=mbAmber!75!black, align=center]
          at (0.45,-2.45) {la droite du modèle,\\ boutons à zéro};
    % --- écarts verticaux ---
    \foreach \mbx/\mby/\mbp/\mbe in {1/2/0/2, 2/5/0/5, 3/8/0/8}{
      \draw[draw=mbRose, line width=1.1pt] (\mbx,\mby) -- (\mbx,\mbp);
      \node[font=\sffamily\scriptsize, text=mbRose, anchor=west]
            at (\mbx+0.08,{(\mby+\mbp)/2}) {\mbe};
    }
    % --- points réels ---
    \foreach \mbx/\mby/\mbnom in {1/2/Léa, 2/5/Sami, 3/8/Nina}{
      \fill[mbBlue] (\mbx,\mby) circle (0.09cm);
      \node[font=\sffamily\scriptsize, text=mbBlue, anchor=north east]
            at (\mbx-0.05,\mby-0.2) {\mbnom};
    }
  \end{scope}
\end{schema}

\begin{aretenir}
Toute la suite du livre repose sur ce couple d'idées : un modèle donne une
réponse, et une erreur mesure de combien cette réponse est fausse.
Apprendre, ce sera toujours et uniquement : \textbf{faire descendre
l'erreur en changeant les paramètres}.
\end{aretenir}

Nous savons mesurer notre bêtise. Reste la question qui compte vraiment :
dans quel sens faut-il tourner les boutons pour que ce nombre baisse ?

\section{De quel côté tourner le bouton ? La pente}

\begin{enclair}
Pour savoir dans quel sens tourner un bouton, il suffit de le pousser un tout
petit peu et de regarder ce que devient l'erreur. Si elle monte, on tournait
du mauvais côté. Si elle descend, on continue.
\end{enclair}

C'est le jeu du « tu chauffes, tu refroidis ». On cache un objet dans une
pièce, l'enfant fait un pas, et on lui dit s'il se rapproche ou s'il
s'éloigne. Avec cette seule information, répétée assez souvent, il finit
toujours par trouver.

Faisons-le pour de vrai sur notre machine à bonbons. Nos boutons sont à
$a = 0$ et $b = 0$, et l'erreur vaut $31$. Poussons le bouton $a$ d'un tout
petit peu, disons de $0{,}01$, sans toucher à $b$.

\begin{exemple}[Pousser un bouton et regarder]
Avec $a = 0{,}01$ et $b = 0$, les trois prédictions deviennent $0{,}01$,
$0{,}02$ et $0{,}03$. Les écarts valent donc $-1{,}99$, $-4{,}98$ et
$-7{,}97$. Leurs carrés valent $3{,}9601$, $24{,}8004$ et $63{,}5209$, dont
la somme fait $92{,}2814$. La nouvelle erreur vaut
\[
  \mathcal{L} \;=\; \frac{92{,}2814}{3} \;\approx\; 30{,}7605 .
\]
Elle est passée de $31$ à $30{,}7605$ : elle a \textbf{baissé} de
$0{,}2395$. Nous poussions donc dans le bon sens. Pour savoir à quelle
vitesse elle baisse, on divise la baisse par la taille de la poussée :
\[
  \frac{-0{,}2395}{0{,}01} \;=\; -23{,}95 .
\]
Ce nombre, $-23{,}95$, dit deux choses à la fois. Son \emph{signe} négatif
dit : « en augmentant $a$, l'erreur diminue ». Sa \emph{taille}, presque
$24$, dit : « et elle diminue vite ».
\end{exemple}

\begin{definition}[La pente][pente]
La \emph{pente} par rapport à un bouton est le nombre qui dit de combien
l'erreur change quand on pousse ce bouton d'une unité. Une pente négative
veut dire que l'erreur descend quand on augmente le bouton. Une pente
positive veut dire qu'elle monte. Une pente proche de zéro veut dire que ce
bouton ne change presque rien.
\end{definition}

\begin{definition}[Le gradient][gradient]
Le \emph{gradient} est la liste de toutes les pentes, une par bouton. S'il y
a deux boutons, le gradient contient deux nombres ; s'il y en a deux
milliards, il en contient deux milliards. Le gradient est donc une sorte de
boussole : il indique, pour chaque bouton à la fois, dans quel sens tourner.
\end{definition}

Pousser chaque bouton à la main pour mesurer sa pente serait épouvantable :
un modèle comme V-JEPA~2.1 a deux milliards de boutons, il faudrait deux
milliards de petites poussées à chaque étape. Heureusement, les
mathématiques savent calculer ces pentes directement, d'un coup, par une
formule. Pour notre droite, les deux formules sont les suivantes :

\begin{notation}[Les deux pentes de notre modèle][pentes-droite]
En notant $e_i = \pred{y}_i - y_i$ l'écart sur l'exemple $i$ :
\[
  \text{pente selon } a \;=\; \frac{2}{N}\sum_i e_i \times x_i,
  \qquad
  \text{pente selon } b \;=\; \frac{2}{N}\sum_i e_i .
\]
\end{notation}

Ces deux formules ont l'air de tomber du ciel. Elles ne sont pourtant pas
mystérieuses, et il vaut la peine de comprendre ce qu'elles racontent.

\begin{intuition}
Regardez la pente selon $b$ : c'est simplement la moyenne des écarts,
multipliée par deux. C'est logique, car le bouton $b$ déplace toute la
droite vers le haut ou vers le bas, de la même quantité pour tout le monde.
Si en moyenne la machine annonce trop de bonbons, il faut baisser $b$.

La pente selon $a$ est la même chose, mais chaque écart y est d'abord
multiplié par $x_i$, le nombre de pièces. C'est l'effet du levier. Le
bouton $a$ fait pivoter la droite autour du point de départ : un enfant
placé loin, comme Nina avec ses trois pièces, monte ou descend trois fois
plus qu'un enfant placé près, comme Léa. Son écart pèse donc trois fois plus
lourd dans la décision.
\end{intuition}

\begin{exemple}[La formule donne-t-elle le même résultat que la poussée ?]
Vérifions. Avec $a = 0$ et $b = 0$, les écarts valent $-2$, $-5$ et $-8$, et
les pièces valent $1$, $2$ et $3$. La pente selon $a$ vaut donc
\begin{align*}
  \frac{2}{3}\bigl[(-2)\times 1 + (-5)\times 2 + (-8)\times 3\bigr]
    &= \frac{2}{3}\bigl[-2 - 10 - 24\bigr] \\[2pt]
    &= \frac{2}{3}\times(-36) = -24 .
\end{align*}
La formule annonce $-24$, et notre petite poussée avait mesuré $-23{,}95$.
Les deux méthodes s'accordent. Le minuscule écart vient de ce que notre
poussée de $0{,}01$ n'était pas infiniment petite ; en poussant encore
moins fort, on collerait encore mieux au $-24$.

Calculons aussi la pente selon $b$ :
\[
  \frac{2}{3}\bigl[(-2) + (-5) + (-8)\bigr]
  = \frac{2}{3}\times(-15) = -10 .
\]
\end{exemple}

Les deux pentes sont négatives. La boussole nous dit donc : augmente $a$,
et augmente $b$ aussi. Voyons maintenant comment transformer ce conseil en
mouvement.

\section{La descente de gradient, calculée à la main}

\begin{enclair}
La règle tient en une phrase : on tourne chaque bouton dans le sens opposé
à sa pente, d'une petite quantité. Puis on recommence.
\end{enclair}

Pourquoi le sens \emph{opposé} ? Parce qu'une pente négative signifie « en
augmentant, l'erreur descend ». Prendre l'opposé du négatif donne du
positif : on augmente. Et une pente positive signifie « en augmentant,
l'erreur monte » ; son opposé nous fait diminuer. Dans les deux cas, on va
vers le bas. C'est la même règle qu'une bille lâchée sur une colline : elle
part toujours dans le sens contraire de la montée.

\begin{definition}[La descente de gradient][descente]
La \emph{descente de gradient} est la méthode qui consiste à répéter
l'opération suivante : calculer l'erreur, calculer le gradient, puis
déplacer chaque paramètre d'un petit pas dans le sens opposé à sa pente.
Pour un paramètre $p$ quelconque, la mise à jour s'écrit
\[
  p_{\text{nouveau}} \;=\; p_{\text{ancien}}
       \;-\; \pas \times \text{pente selon } p ,
\]
où $\pas$ est un petit nombre positif choisi par nous, qui règle la
longueur du pas.
\end{definition}

La figure~\ref{fig:vallee} montre cette idée sous sa forme la plus parlante :
l'erreur dessine une vallée, et le modèle est une bille qui la descend.

\begin{schema}{L'erreur vue comme une vallée. En allant vers la droite ou
vers la gauche, on change la valeur d'un bouton ; la hauteur de la courbe
donne l'erreur correspondante. La bille est notre modèle. La droite en
pointillés est la pente à l'endroit où se trouve la bille : elle monte vers
la droite, donc la pente est positive, donc il faut aller à gauche. Le
meilleur réglage est tout en bas de la vallée, là où la courbe est
plate.}[vallee]
  \begin{scope}[x=1.15cm, y=0.72cm]
    % --- la vallée ---
    \draw[draw=mbSlate!55, line width=1.2pt, smooth, samples=80,
          domain=-2.9:2.9] plot (\x, {0.5*\x*\x});
    % --- axes discrets ---
    \draw[->, draw=mbSlate!70, line width=0.6pt] (-3.3,0) -- (3.4,0)
          node[right, font=\sffamily\scriptsize, text=mbSlate]
          {valeur du bouton};
    \draw[->, draw=mbSlate!70, line width=0.6pt] (-3.3,-0.3) -- (-3.3,4.6)
          node[above, font=\sffamily\scriptsize, text=mbSlate] {erreur};
    % --- tangente au point (2 ; 2), de pente 2 ---
    \draw[draw=mbAmber!85!black, line width=1pt, dashed]
          (1.25,0.5) -- (2.75,3.5);
    \node[font=\sffamily\scriptsize, text=mbAmber!75!black, anchor=west]
          at (2.8,3.4) {la pente ici};
    % --- la bille ---
    \fill[mbBlue] (2,2) circle (0.13cm);
    \node[font=\sffamily\scriptsize, text=mbBlue, anchor=south west]
          at (2.12,2.12) {le modèle};
    % --- sens du déplacement (en vert, pour ne pas le confondre avec
    %     la tangente, qui est ambre et en pointillés) ---
    \draw[->, draw=mbTeal, line width=1.1pt] (1.85,1.55) -- (1.05,1.05);
    \node[font=\sffamily\scriptsize\itshape, text=mbTeal,
          anchor=east] at (1.0,1.35) {on va par là};
    % --- le fond ---
    \fill[mbGreen] (0,0) circle (0.1cm);
    \node[font=\sffamily\scriptsize, text=mbGreen!60!black, align=center]
          at (0,-0.85) {le meilleur\\ réglage};
  \end{scope}
\end{schema}

Nous avons tout ce qu'il faut. Choisissons une longueur de pas de $0{,}1$ et
faisons deux étapes complètes, en écrivant chaque calcul.

\begin{exemple}[Première étape, tout en détail]
Nous partons de $a = 0$ et $b = 0$. Nous avons déjà tout calculé :
l'erreur vaut $31$, la pente selon $a$ vaut $-24$, la pente selon $b$ vaut
$-10$. Appliquons la règle avec un pas de $0{,}1$ :
\[
  a_{\text{nouveau}} = 0 - 0{,}1 \times (-24) = 0 + 2{,}4 = 2{,}4 ,
\]
\[
  b_{\text{nouveau}} = 0 - 0{,}1 \times (-10) = 0 + 1 = 1 .
\]
Nos deux boutons sont maintenant réglés à $a = 2{,}4$ et $b = 1$.
\end{exemple}

\begin{exemple}[A-t-on vraiment progressé ?]
Vérifions en recalculant l'erreur avec les nouveaux boutons. Les trois
prédictions valent :
\[
  2{,}4 \times 1 + 1 = 3{,}4 ; \quad
  2{,}4 \times 2 + 1 = 5{,}8 ; \quad
  2{,}4 \times 3 + 1 = 8{,}2 .
\]
Les vraies valeurs étaient $2$, $5$ et $8$. Les écarts valent donc $1{,}4$,
$0{,}8$ et $0{,}2$, et leurs carrés $1{,}96$, $0{,}64$ et $0{,}04$. La
nouvelle erreur vaut
\[
  \mathcal{L} = \frac{1{,}96 + 0{,}64 + 0{,}04}{3}
              = \frac{2{,}64}{3} = 0{,}88 .
\]
Nous sommes passés de $31$ à $0{,}88$ en une seule étape. L'erreur a été
divisée par plus de trente-cinq. La machine, qui annonçait zéro bonbon à
tout le monde, en annonce maintenant trois et demi à Léa et huit à Nina :
elle a déjà presque compris.
\end{exemple}

\begin{exemple}[Deuxième étape]
Recommençons exactement le même travail depuis $a = 2{,}4$ et $b = 1$. Les
écarts, que nous venons de calculer, valent $1{,}4$, $0{,}8$ et $0{,}2$. Ils
sont maintenant positifs : la machine annonce trop de bonbons.

Pente selon $a$ :
\[
  \frac{2}{3}\bigl[1{,}4 \times 1 + 0{,}8 \times 2 + 0{,}2 \times 3\bigr]
  = \frac{2}{3}\bigl[1{,}4 + 1{,}6 + 0{,}6\bigr]
  = \frac{2}{3} \times 3{,}6 = 2{,}4 .
\]
Pente selon $b$ :
\[
  \frac{2}{3}\bigl[1{,}4 + 0{,}8 + 0{,}2\bigr]
  = \frac{2}{3} \times 2{,}4 = 1{,}6 .
\]
Les deux pentes sont cette fois positives : il faut donc \emph{baisser} les
deux boutons.
\[
  a_{\text{nouveau}} = 2{,}4 - 0{,}1 \times 2{,}4 = 2{,}4 - 0{,}24 = 2{,}16 ,
\]
\[
  b_{\text{nouveau}} = 1 - 0{,}1 \times 1{,}6 = 1 - 0{,}16 = 0{,}84 .
\]
Et la nouvelle erreur vaut $0{,}496$, contre $0{,}88$ juste avant : elle a
encore baissé.
\end{exemple}

Récapitulons les deux étapes dans un seul tableau, que vous pourrez
comparer tout à l'heure à la sortie du programme.

\begin{center}
\begin{tabular}{cccccc}
\toprule
\textbf{Étape} & $a$ & $b$ & \textbf{erreur} &
\textbf{pente selon } $a$ & \textbf{pente selon } $b$ \\
\midrule
0 & $0$      & $0$      & $31$      & $-24$   & $-10$  \\
1 & $2{,}4$  & $1$      & $0{,}88$  & $2{,}4$ & $1{,}6$ \\
2 & $2{,}16$ & $0{,}84$ & $0{,}496$ & $-0{,}48$ & $0{,}32$ \\
\bottomrule
\end{tabular}
\end{center}

\begin{piege}
Regardez la colonne du bouton $a$ : il monte à $2{,}4$, puis \emph{redescend}
à $2{,}16$, alors que la bonne réponse est $3$. On dirait qu'il s'éloigne !
Ne vous inquiétez pas, et surtout ne jugez jamais un entraînement en
surveillant un paramètre isolé. Les boutons bougent ensemble, ils se
gênent et se rattrapent l'un l'autre. Le seul chiffre qui compte est
l'erreur, et elle, elle descend sans arrêt : $31$, puis $0{,}88$, puis
$0{,}496$. Au bout de deux cents étapes, $a$ vaut $2{,}993$ et $b$ vaut
$-0{,}985$ : la machine a bien retrouvé le $3$ et le $-1$ que nous
cherchions.
\end{piege}

\begin{enia}
Ce que nous venons de faire à la main sur deux boutons est \emph{exactement}
ce que fait l'entraînement de V-JEPA~2.1, sans un seul changement d'idée.
Les seules différences sont des différences de taille : deux milliards de
boutons au lieu de deux, une erreur calculée sur des morceaux de vidéo au
lieu de trois enfants, et cent trente-cinq mille étapes au lieu de deux.
La boucle, elle, est identique : je prédis, je mesure mon erreur, je calcule
mes pentes, je fais un petit pas, je recommence.
\end{enia}

Il nous reste à parler du seul réglage que nous avons choisi arbitrairement,
et qui est pourtant le plus délicat de tous : la longueur du pas.

\section{La taille des pas : le taux d'apprentissage}

\begin{definition}[Le taux d'apprentissage][taux]
Le \emph{taux d'apprentissage}, noté $\pas$, est le petit nombre par lequel
on multiplie la pente avant de bouger un paramètre. Il décide de la longueur
des pas. C'est nous qui le choisissons : la machine ne l'apprend pas.
\end{definition}

\begin{intuition}
Vous descendez une colline dans le brouillard. Vous ne voyez que le sol sous
vos pieds, donc vous ne connaissez que la pente à l'endroit où vous êtes.
Avec des pas minuscules, vous descendrez, mais la nuit tombera avant que
vous arriviez. Avec des pas de sept lieues, vous sauterez par-dessus la
vallée et vous vous retrouverez plus haut qu'avant, sur l'autre versant. Il
existe une bonne taille de pas entre les deux, et la trouver fait partie du
métier.
\end{intuition}

La figure~\ref{fig:taille-des-pas} met les trois situations côte à côte.

\begin{schema}{Trois façons de descendre la même vallée. À gauche, les pas
sont trop courts : la bille avance à peine et il faudra une éternité. Au
milieu, les pas sont bien choisis : la bille rejoint le fond en quelques
sauts. À droite, les pas sont trop longs : la bille saute par-dessus le fond,
remonte plus haut de l'autre côté, et s'échappe de la vallée au lieu d'y
descendre.}[taille-des-pas]
  \begin{scope}[x=0.42cm, y=0.30cm]
    % ================= panneau 1 : pas trop petit =================
    % depart en x = 3, chaque pas ne retire que 12 % du chemin restant
    \begin{scope}
      \draw[draw=mbSlate!45, line width=1pt, smooth, samples=60,
            domain=-3.7:3.7] plot (\x, {0.5*\x*\x});
      \foreach \mbx/\mby in {3.0/4.5, 2.64/3.485, 2.323/2.698,
                             2.044/2.089, 1.799/1.618}{
        \fill[mbBlue] (\mbx,\mby) circle (0.07cm);
      }
      \draw[fleche accent] (3.05,5.35) -- (2.0,5.35);
      \node[font=\sffamily\scriptsize\bfseries, text=mbSlate]
            at (0,-1.6) {pas trop petit};
      \node[font=\sffamily\scriptsize, text=mbSlate!85, align=center]
            at (0,-3.4) {on descend,\\ mais si lentement};
    \end{scope}
    % ================= panneau 2 : pas bien choisi =================
    % meme depart, mais chaque pas retire la moitie du chemin restant
    \begin{scope}[xshift=4.0cm]
      \draw[draw=mbSlate!45, line width=1pt, smooth, samples=60,
            domain=-3.7:3.7] plot (\x, {0.5*\x*\x});
      \foreach \mbx/\mby in {3.0/4.5, 1.5/1.125, 0.75/0.281,
                             0.375/0.070, 0.188/0.018}{
        \fill[mbGreen] (\mbx,\mby) circle (0.07cm);
      }
      \draw[fleche accent] (3.05,5.35) -- (0.55,5.35);
      \node[font=\sffamily\scriptsize\bfseries, text=mbGreen!55!black]
            at (0,-1.6) {pas bien choisi};
      \node[font=\sffamily\scriptsize, text=mbSlate!85, align=center]
            at (0,-3.4) {on atteint le fond\\ en quelques pas};
    \end{scope}
    % ================= panneau 3 : pas trop grand =================
    % chaque pas depasse le fond et repart plus haut de l'autre cote
    \begin{scope}[xshift=8.0cm]
      \draw[draw=mbSlate!45, line width=1pt, smooth, samples=60,
            domain=-3.7:3.7] plot (\x, {0.5*\x*\x});
      \foreach \mbx/\mby in {1.5/1.125, -1.95/1.901, 2.535/3.213,
                             -3.296/5.434}{
        \fill[mbRose] (\mbx,\mby) circle (0.07cm);
      }
      \draw[fleche accent, draw=mbRose] (1.28,1.30) -- (-1.68,1.76);
      \draw[fleche accent, draw=mbRose] (-1.70,2.16) -- (2.28,3.09);
      \draw[fleche accent, draw=mbRose] (2.30,3.52) -- (-3.02,5.29);
      \node[font=\sffamily\scriptsize\bfseries, text=mbRose]
            at (0,-1.6) {pas trop grand};
      \node[font=\sffamily\scriptsize, text=mbSlate!85, align=center]
            at (0,-3.4) {on saute par-dessus\\ et l'on remonte};
    \end{scope}
  \end{scope}
\end{schema}

Ce ne sont pas des dessins pour faire joli : refaisons une seule étape de
notre machine à bonbons avec trois pas différents, et regardons ce que
devient l'erreur.

\begin{exemple}[Trois pas, trois destins]
Nous partons toujours de $a = 0$, $b = 0$, avec une erreur de $31$ et des
pentes de $-24$ et $-10$. Une seule étape suffit à tout révéler.

\begin{center}
\begin{tabular}{cccc}
\toprule
\textbf{Pas} $\pas$ & \textbf{nouveau } $a$ & \textbf{nouveau } $b$ &
\textbf{erreur après une étape} \\
\midrule
$0{,}001$ & $0{,}024$ & $0{,}01$ & $30{,}33$ \\
$0{,}1$   & $2{,}4$   & $1$      & $\bm{0{,}88}$ \\
$1$       & $24$      & $10$     & $3\,103$ \\
\bottomrule
\end{tabular}
\end{center}

Avec un pas de $0{,}001$, l'erreur passe de $31$ à $30{,}33$ : elle a bougé
d'à peine deux pour cent. À ce rythme, il faudrait des milliers d'étapes.
Avec un pas de $0{,}1$, elle tombe à $0{,}88$ : c'est le bon réglage. Avec
un pas de $1$, elle grimpe à $3\,103$, soit cent fois pire qu'au départ ! Et
si l'on s'obstine, l'étape suivante donne une erreur de $316\,063$. La
machine ne se contente pas d'échouer : elle s'envole.
\end{exemple}

\begin{piege}
Un taux d'apprentissage trop grand est la panne numéro un des débutants, et
elle est facile à reconnaître : au lieu de descendre doucement, l'erreur
grossit d'étape en étape, jusqu'à afficher des nombres absurdes, voire le
mot \co{python}{nan}, qui veut dire « ce n'est plus un nombre du tout ».
Le réflexe correct, dans ce cas, est presque toujours le même : diviser le
taux d'apprentissage par dix et relancer.
\end{piege}

\begin{surlesdonnees}
V-JEPA~2.1 ne garde pas le même pas du début à la fin. Pendant les douze
mille premières étapes, il l'augmente doucement de $10^{-4}$ à
$5{,}25 \times 10^{-4}$ : on appelle cela une montée en douceur, et cela
évite de tout casser au tout début, quand les boutons sont encore réglés
n'importe comment. Ensuite il le laisse constant très longtemps. Enfin, sur les douze
mille dernières étapes, il le fait redescendre jusqu'à $10^{-6}$, c'est-à-dire
presque zéro : les pas deviennent minuscules et le modèle se pose tout
doucement au fond de la vallée, au lieu de continuer à sautiller autour.
Cette dernière phase porte un nom que nous retrouverons souvent : la phase
de refroidissement.
\end{surlesdonnees}

Nous savons faire un pas, et le faire de la bonne longueur. Il reste à
savoir sur combien de données on le calcule, et combien de fois on
recommence.

\section{Le rythme de l'entraînement : le lot et l'époque}

Jusqu'ici, nous avions trois enfants, et nous avons calculé l'erreur sur les
trois à chaque étape. Avec trois exemples, c'est facile. Avec cent
quarante-deux millions d'images, c'est impossible : regarder toute la
collection avant de bouger un seul bouton prendrait des jours pour une seule
étape.

\begin{enclair}
La solution est simple et un peu paresseuse : au lieu de regarder toutes les
données, on en prend une petite poignée au hasard, on calcule l'erreur sur
cette poignée-là, et on fait quand même un pas. La direction est un peu
approximative, mais on avance beaucoup plus souvent — et au bout du compte,
on va bien plus vite.
\end{enclair}

\begin{definition}[Le lot][lot]
Un \emph{lot} est la petite poignée d'exemples que l'on regarde ensemble
pour calculer une erreur et faire un pas. En anglais on dit \emph{batch}, et
ce mot revient partout dans le code. Le nombre d'exemples dans la poignée
s'appelle la \emph{taille de lot}.
\end{definition}

\begin{definition}[Le pas, ou itération][iteration]
Un \emph{pas} — on dit aussi une \emph{itération} — est un tour complet de
la boucle : on prend un lot, on prédit, on mesure l'erreur, on calcule les
pentes, on bouge les paramètres. Un pas, un lot.
\end{definition}

\begin{piege}
Le mot « pas » sert malheureusement à deux choses, et il faut les séparer
dans votre tête une bonne fois pour toutes. \emph{Faire un pas}, c'est faire
un tour de boucle : c'est ce que nous venons de définir, et cela se compte
($135\,000$ pas pour V-JEPA~2.1). \emph{La longueur du pas}, c'est le taux
d'apprentissage $\pas$ de la section précédente : c'est de combien on bouge
les boutons pendant ce tour, et cela se règle. On peut donc faire beaucoup
de pas très courts, ou peu de pas très longs. Quand un texte parle de
« pas », demandez-vous toujours s'il compte des tours ou s'il mesure une
longueur.
\end{piege}

\begin{definition}[L'époque][epoque]
Une \emph{époque} est le temps qu'il faut pour que le modèle ait vu une fois
chacune des données disponibles. Si la collection contient mille images et
que les lots en contiennent cent, alors une époque vaut dix pas.
\end{definition}

\begin{schema}{Toutes les données de la collection, découpées en lots. Le
modèle en prend un, fait un pas, en prend un autre, refait un pas. Quand il
a épuisé tous les lots, il a vu chaque donnée une fois : cela s'appelle une
époque, et il peut alors recommencer un nouveau tour.}[lot-epoque]
  \begin{scope}[x=1cm, y=1cm]
    \foreach \mbg in {0,1,2}{
      \begin{scope}[xshift=\mbg*3.55cm]
        \foreach \mbk in {0,1,2,3}{
          \draw[draw=mbBlue!70, fill=mbBlue!14, line width=0.7pt,
                rounded corners=1.5pt]
                (\mbk*0.78,0) rectangle ++(0.62,0.62);
        }
        \draw[accolade bas] (0,-0.12) -- (3.02,-0.12);
        \node[font=\sffamily\scriptsize, text=mbSlate] at (1.51,-0.66)
             {lot n\textsuperscript{o}\,\the\numexpr\mbg+1\relax};
        \draw[fleche accent] (1.51,0.95) -- (1.51,0.72);
        \node[font=\sffamily\scriptsize\itshape, text=mbAmber!75!black]
              at (1.51,1.22) {un pas};
      \end{scope}
    }
    \draw[accolade bas] (0,-1.05) -- (9.62,-1.05);
    \node[font=\sffamily\small\bfseries, text=mbGreen!55!black]
          at (4.81,-1.7) {une époque : chaque donnée a été vue une fois};
  \end{scope}
\end{schema}

\begin{exemple}[Compter des époques sur un petit jeu]
Vous disposez de $1\,000$ photos et vous choisissez des lots de $100$
photos. Chaque pas consomme $100$ photos, il faut donc
\[
  \frac{1\,000}{100} = 10 \text{ pas pour faire une époque.}
\]
Si vous entraînez pendant $50$ époques, votre programme aura effectué
$50 \times 10 = 500$ pas au total, et chaque photo aura été vue $50$ fois.
\end{exemple}

\begin{surlesdonnees}
Faisons le même calcul pour V-JEPA~2.1, et le résultat va sans doute vous
surprendre. Le modèle fait $135\,000$ pas dans sa phase principale, et
chaque lot contient $2\,304$ images. Il a donc regardé
\[
  135\,000 \times 2\,304 = 311\,040\,000 \text{ images en tout,}
\]
soit trois cent onze millions. Or sa collection d'images en compte $142$
millions. Le nombre d'époques vaut donc
\[
  \frac{311\,040\,000}{142\,000\,000} \approx 2{,}2 .
\]
Chaque image n'est vue qu'un peu plus de deux fois dans toute la vie du
modèle. Du côté des vidéos, c'est encore plus frappant : $135\,000$ pas
avec $128$ clips par lot font $17\,280\,000$ clips, pour une collection de
$19$ millions de vidéos. Le modèle ne voit donc même pas une fois chaque
vidéo !

C'est pour cela que les articles de recherche modernes ne parlent presque
plus d'époques, mais comptent en pas. Quand les données sont quasi
inépuisables, repasser deux fois sur la même image n'a plus grand intérêt :
mieux vaut en voir une nouvelle.
\end{surlesdonnees}

\begin{plusloin}
Prendre un lot au hasard plutôt que toutes les données rend le gradient un
peu bruité : la direction indiquée n'est pas tout à fait la bonne, elle
tremble. On appelle cela la descente de gradient \emph{stochastique}, un mot
savant qui veut simplement dire « au hasard ». Ce tremblement, contre toute
attente, aide souvent : il empêche la bille de rester coincée dans un petit
creux sans intérêt et lui fait trouver de meilleures vallées.
\end{plusloin}

Nous savons entraîner. Il reste un dernier danger, et c'est le plus sournois
de tous : une machine peut très bien faire baisser son erreur tout en
n'ayant strictement rien compris.

\section{Apprendre par cœur n'est pas apprendre}

\begin{enclair}
Un élève peut apprendre par cœur les réponses des exercices du livre sans
comprendre la leçon. Tant qu'on l'interroge sur ces exercices-là, il a tout
juste. Le jour où on lui pose une question nouvelle, il s'effondre. Une
machine peut faire exactement pareil.
\end{enclair}

Revenons à nos trois enfants. Notre droite n'a que deux boutons : elle est
trop simple pour tricher, elle est bien obligée de trouver une règle
générale. Mais donnons à un modèle beaucoup, beaucoup de boutons, et il
pourra faire passer une courbe tordue exactement par chacun des trois
points. Son erreur tombera à zéro. Et pourtant, si l'on glisse quatre pièces
dans la machine, ce modèle-là annoncera n'importe quoi.

\begin{definition}[Le sur-apprentissage][surapprentissage]
Le \emph{sur-apprentissage} est ce qui arrive quand un modèle apprend par
cœur les exemples qu'on lui a montrés, y compris leurs petits défauts, au
lieu d'apprendre la règle générale qui se cache derrière. Il est alors
excellent sur ce qu'il a déjà vu, et mauvais sur tout le reste.
\end{definition}

\begin{schema}{Les mêmes points, expliqués de trois façons. À gauche, le
modèle est trop simple : il ne suit même pas la tendance, il passe à côté de
tout le monde. Au milieu, il a saisi la règle générale : il ne touche pas
tous les points, et c'est très bien ainsi. À droite, il passe exactement par
chaque point, son erreur est nulle — mais entre les points il fait n'importe
quoi, et sur un point nouveau il se trompera lourdement.}[surapprentissage]
  \begin{scope}[x=0.76cm, y=0.66cm]
    % ============ panneau 1 : trop simple ============
    \begin{scope}
      \draw[draw=mbSlate!35, line width=0.6pt] (0,0) rectangle (4.4,4.6);
      \draw[draw=mbRose, line width=1.2pt] (0.25,2.45) -- (4.15,2.45);
      \foreach \mbx/\mby in {0.3/0.9, 0.9/1.9, 1.5/1.7, 2.1/2.9,
                             2.7/2.6, 3.3/3.7, 3.9/3.5}{
        \fill[mbBlue] (\mbx,\mby) circle (0.075cm);
      }
      \node[font=\sffamily\scriptsize\bfseries, text=mbRose]
            at (2.2,-0.55) {trop simple};
      \node[font=\sffamily\scriptsize, text=mbSlate!85, align=center]
            at (2.2,-1.4) {il n'a rien\\ appris du tout};
    \end{scope}
    % ============ panneau 2 : juste ============
    \begin{scope}[shift={(5.1,0)}]
      \draw[draw=mbSlate!35, line width=0.6pt] (0,0) rectangle (4.4,4.6);
      \draw[draw=mbGreen, line width=1.2pt] (0.25,0.9) -- (4.15,4.0);
      \foreach \mbx/\mby in {0.3/0.9, 0.9/1.9, 1.5/1.7, 2.1/2.9,
                             2.7/2.6, 3.3/3.7, 3.9/3.5}{
        \fill[mbBlue] (\mbx,\mby) circle (0.075cm);
      }
      \node[font=\sffamily\scriptsize\bfseries, text=mbGreen!55!black]
            at (2.2,-0.55) {bien réglé};
      \node[font=\sffamily\scriptsize, text=mbSlate!85, align=center]
            at (2.2,-1.4) {il a compris\\ la tendance};
    \end{scope}
    % ============ panneau 3 : par cœur ============
    \begin{scope}[shift={(10.2,0)}]
      \draw[draw=mbSlate!35, line width=0.6pt] (0,0) rectangle (4.4,4.6);
      \draw[draw=mbAmber!85!black, line width=1.2pt, smooth]
            plot coordinates {(0.3,0.9) (0.9,1.9) (1.5,1.7) (2.1,2.9)
                              (2.7,2.6) (3.3,3.7) (3.9,3.5)};
      \foreach \mbx/\mby in {0.3/0.9, 0.9/1.9, 1.5/1.7, 2.1/2.9,
                             2.7/2.6, 3.3/3.7, 3.9/3.5}{
        \fill[mbBlue] (\mbx,\mby) circle (0.075cm);
      }
      \node[font=\sffamily\scriptsize\bfseries, text=mbAmber!80!black]
            at (2.2,-0.55) {appris par cœur};
      \node[font=\sffamily\scriptsize, text=mbSlate!85, align=center]
            at (2.2,-1.4) {erreur nulle,\\ mais rien compris};
    \end{scope}
  \end{scope}
\end{schema}

Comment se protéger de ce piège ? En trichant le moins possible : on cache
au modèle une partie des données, et on ne les lui montre jamais pendant
l'entraînement. Elles serviront uniquement à le juger.

\begin{definition}[Jeu d'entraînement et jeu de validation][jeux]
Le \emph{jeu d'entraînement} rassemble les données que le modèle a le droit
de voir pour régler ses boutons. Le \emph{jeu de validation}, parfois appelé
jeu de test, rassemble des données qu'il ne voit jamais pendant
l'entraînement, et qui servent uniquement à mesurer honnêtement ce qu'il
vaut. C'est la différence entre les exercices du cahier et l'examen de fin
d'année.
\end{definition}

\begin{piege}
La faute la plus grave, dans tout l'apprentissage automatique, est de régler
quoi que ce soit en regardant le jeu de validation. Si vous essayez trente
taux d'apprentissage et que vous gardez celui qui donne le meilleur score de
validation, alors ce score n'est plus honnête : vous avez, sans le vouloir,
appris par cœur votre examen. Les vrais résultats se mesurent toujours sur
des données que rien ni personne n'a servi à choisir.
\end{piege}

\begin{enia}
L'article de V-JEPA~2.1 prend cette précaution très au sérieux, et cela se
voit dans ses annexes. Pour le suivi d'objets dans les vidéos, par exemple,
les auteurs cherchent les meilleurs réglages sur la partie entraînement du
jeu DAVIS, puis appliquent tels quels ces réglages à tous les jeux
d'évaluation, sans plus jamais y toucher. C'est un détail d'annexe, et
pourtant c'est ce détail qui rend leurs chiffres croyables.
\end{enia}

Nous avons maintenant tout le vocabulaire. Le moment est venu de vérifier
que nos calculs de papier tiennent devant une vraie machine.

\section{Passons au code}

Nous allons refaire, en Python, les deux étapes que nous venons de calculer
à la main. Nous utiliserons \co{python}{torch}, la bibliothèque PyTorch, qui
sait faire deux choses très utiles : manipuler des tenseurs, et surtout
calculer les pentes toute seule. Nous n'aurons donc plus jamais à écrire les
formules du gradient nous-mêmes.

Commençons par le plus important : vérifier que la petite poussée du début
du chapitre donne bien ce que nous avions trouvé.

\begin{code}{python}[Mesurer une pente en poussant un bouton]
import torch

# Les trois enfants : pieces glissees et bonbons recus.
x = torch.tensor([1.0, 2.0, 3.0])
y = torch.tensor([2.0, 5.0, 8.0])

def erreur(a, b):
    prediction = a * x + b
    return ((prediction - y) ** 2).mean().item()

avant = erreur(0.00, 0.0)   # les boutons a zero
apres = erreur(0.01, 0.0)   # on pousse le bouton a de 0,01

print(f"erreur avant la poussee : {avant:.4f}")   # -> 31.0000
print(f"erreur apres la poussee : {apres:.4f}")   # -> 30.7605
print(f"pente mesuree           : {(apres - avant) / 0.01:.2f}")  # -> -23.95
\end{code}

Les trois nombres affichés sont exactement ceux de notre exemple :
$31$, $30{,}7605$ et $-23{,}95$. Le calcul à la main et la machine sont
d'accord.

Voici maintenant la boucle d'entraînement complète. Elle tient en une
quinzaine de lignes, et c'est littéralement la même boucle que celle qui
entraîne V-JEPA~2.1 : seules les tailles changent.

\begin{code}{python}[Deux étapes de descente de gradient]
import torch

x = torch.tensor([1.0, 2.0, 3.0])
y = torch.tensor([2.0, 5.0, 8.0])

# Nos deux boutons. « requires_grad » dit a PyTorch :
# surveille ce nombre, je te demanderai sa pente.
a = torch.tensor(0.0, requires_grad=True)
b = torch.tensor(0.0, requires_grad=True)
pas = 0.1

# 1. je prends mes donnees : ici, les trois enfants d'un seul coup.
for etape in range(2):
    prediction = a * x + b                       # 2. je calcule ma prediction
    erreur = ((prediction - y) ** 2).mean()      # 3. je mesure mon erreur
    erreur.backward()                            # 4. je calcule les pentes

    print(f"etape {etape} : a = {a.item():.2f}, b = {b.item():.2f}, "
          f"erreur = {erreur.item():.4f}, "
          f"pente_a = {a.grad.item():.2f}, pente_b = {b.grad.item():.2f}")

    with torch.no_grad():                        # 5. je bouge les boutons
        a -= pas * a.grad
        b -= pas * b.grad

    a.grad.zero_()          # puis j'efface les pentes pour le tour suivant
    b.grad.zero_()

print(f"apres 2 etapes : a = {a.item():.2f}, b = {b.item():.2f}")
\end{code}

Ce programme affiche exactement ceci :

\begin{sortie}
etape 0 : a = 0.00, b = 0.00, erreur = 31.0000, pente_a = -24.00, pente_b = -10.00
etape 1 : a = 2.40, b = 1.00, erreur = 0.8800, pente_a = 2.40, pente_b = 1.60
apres 2 etapes : a = 2.16, b = 0.84
\end{sortie}

Comparez ligne à ligne avec le tableau de la section précédente : le $31$,
le $-24$, le $-10$, le $2{,}4$, le $0{,}88$, le $2{,}16$ et le $0{,}84$ y
sont tous. Rien n'a été caché, et la machine n'a rien fait de plus que ce
que nous avons fait au crayon.

\begin{piege}
La ligne \co{python}{a.grad.zero_()} paraît anodine, et pourtant l'oublier
est l'une des erreurs les plus courantes chez les débutants. PyTorch
\emph{additionne} les pentes qu'il calcule au lieu de les remplacer. Si vous
n'effacez pas, l'étape suivante utilisera la somme de deux pentes, puis de
trois, et vos pas grandiront tout seuls jusqu'à faire exploser
l'entraînement. Le symptôme est le même que celui d'un taux d'apprentissage
trop grand, ce qui rend la panne difficile à diagnostiquer.
\end{piege}

Laissons enfin tourner la même boucle bien plus longtemps, pour voir la
machine retrouver la règle des bonbons.

\begin{code}{python}[Deux cents étapes, et la règle apparaît]
import torch

x = torch.tensor([1.0, 2.0, 3.0])
y = torch.tensor([2.0, 5.0, 8.0])
a = torch.tensor(0.0, requires_grad=True)
b = torch.tensor(0.0, requires_grad=True)

for _ in range(200):
    erreur = ((a * x + b - y) ** 2).mean()
    erreur.backward()
    with torch.no_grad():
        a -= 0.1 * a.grad
        b -= 0.1 * b.grad
    a.grad.zero_()
    b.grad.zero_()

print(f"a = {a.item():.3f}   b = {b.item():.3f}")   # -> a = 2.993   b = -0.985
print(f"erreur = {erreur.item():.6f}")              # -> erreur = 0.000034
\end{code}

Le bouton $a$ vaut $2{,}993$ et le bouton $b$ vaut $-0{,}985$. La machine a
retrouvé, toute seule, la règle que nous avions devinée de tête au début du
chapitre : trois bonbons par pièce, moins un pour la machine. Personne ne
lui a jamais soufflé la réponse. Elle a seulement descendu la vallée, deux
cents fois de suite.

\begin{aretenir}
Retenez cette boucle en cinq gestes, car c'est la seule qu'il y ait dans
tout l'apprentissage automatique, du plus petit exemple au plus gros modèle
du monde :

\begin{enumerate}
  \item je prends un lot de données ;
  \item je calcule ma prédiction ;
  \item je mesure mon erreur ;
  \item je calcule les pentes ;
  \item je bouge chaque bouton d'un petit pas dans le sens opposé à sa pente.
\end{enumerate}

Puis je recommence. Cent trente-cinq mille fois, dans le cas de V-JEPA~2.1.
\end{aretenir}

Il manque pourtant quelque chose d'énorme à tout ce que nous venons de
faire, et vous l'avez peut-être déjà remarqué. À chaque étape, nous avions
besoin de connaître la bonne réponse : les deux bonbons de Léa, les cinq de
Sami, les huit de Nina. Sans elles, impossible de calculer le moindre écart,
donc impossible de calculer la moindre erreur, donc impossible d'apprendre
quoi que ce soit.

Mais qui va écrire la bonne réponse à côté d'un milliard de vidéos ?
Personne. Aucune équipe au monde, aucun budget, aucune quantité de patience
ne suffirait. C'est précisément le mur contre lequel s'arrête tout ce que
nous venons d'apprendre — et c'est le mur que le chapitre suivant va
traverser.

\phantomsection
\section*{Exercices}
\addcontentsline{toc}{section}{Exercices}

\begin{exercice}[Calculer une erreur]
Un modèle prédit $\pred{y} = 4$, $\pred{y} = 7$ et $\pred{y} = 9$, alors que
les vraies valeurs sont $3$, $9$ et $9$. Quelle est l'erreur, au sens de la
moyenne des carrés des écarts ?
\end{exercice}

\begin{solution}
Les écarts valent $4 - 3 = 1$, puis $7 - 9 = -2$, puis $9 - 9 = 0$. Leurs
carrés valent $1$, $4$ et $0$, dont la somme fait $5$. Comme il y a trois
exemples,
\[
  \mathcal{L} = \frac{5}{3} \approx 1{,}67 .
\]
Remarquez que le deuxième écart, bien que négatif, contribue positivement à
l'erreur : c'est tout l'intérêt du carré.
\end{solution}

\begin{exercice}[Faire un pas]
Un paramètre vaut actuellement $p = 5$. On calcule sa pente et l'on trouve
$+3$. Le taux d'apprentissage vaut $0{,}2$. Que vaut le paramètre après la
mise à jour ? A-t-il augmenté ou diminué, et pourquoi était-ce le bon sens ?
\end{exercice}

\begin{solution}
On applique la règle $p_{\text{nouveau}} = p_{\text{ancien}} - \pas \times
\text{pente}$, soit
\[
  p_{\text{nouveau}} = 5 - 0{,}2 \times 3 = 5 - 0{,}6 = 4{,}4 .
\]
Le paramètre a diminué. C'était le bon sens, car une pente positive signifie
« si j'augmente ce bouton, l'erreur monte ». Pour faire descendre l'erreur,
il fallait donc baisser le bouton.
\end{solution}

\begin{exercice}[Compter des époques]
Un jeu de données contient $50\,000$ images et l'on choisit des lots de
$250$ images. Combien de pas faut-il pour une époque ? Et si l'on entraîne
pendant $30\,000$ pas, combien d'époques cela représente-t-il ?
\end{exercice}

\begin{solution}
Une époque demande $50\,000 / 250 = 200$ pas. Avec $30\,000$ pas au total,
on effectue
\[
  \frac{30\,000}{200} = 150 \text{ époques,}
\]
et chaque image aura donc été vue $150$ fois. C'est un régime très différent
de celui de V-JEPA~2.1, qui ne dépasse pas deux époques et demie : quand les
données sont peu nombreuses, on repasse dessus souvent, et le risque de
sur-apprentissage augmente d'autant.
\end{solution}

\begin{exercice}[Diagnostiquer une panne]
Vous lancez un entraînement et vous voyez défiler ces erreurs :
$18{,}4$ puis $91{,}7$ puis $2\,340$ puis \co{python}{nan}. Que se passe-t-il
très probablement, et quel est le premier réglage à corriger ?
\end{exercice}

\begin{solution}
L'erreur ne descend pas, elle explose : le modèle saute par-dessus le fond
de la vallée et remonte de plus en plus haut à chaque étape, jusqu'à sortir
des nombres représentables par la machine, ce qui produit le
\co{python}{nan}. C'est la signature d'un taux d'apprentissage trop grand.
Le premier geste est de le diviser par dix et de relancer ; si le problème
persiste, on divise encore par dix. Attention toutefois : oublier
d'effacer les pentes avec \co{python}{zero_()} produit exactement le même
symptôme, et vaut la peine d'être vérifié dans la foulée.
\end{solution}

\begin{exercice}[Le piège de l'examen]
Vous disposez d'un jeu d'entraînement et d'un jeu de validation. Vous
essayez cinquante taux d'apprentissage différents et vous annoncez, comme
résultat final de votre travail, le meilleur score obtenu sur le jeu de
validation. Pourquoi ce chiffre est-il trompeur, et que fallait-il faire ?
\end{exercice}

\begin{solution}
En choisissant le réglage qui donne le meilleur score de validation, vous
avez utilisé le jeu de validation pour prendre une décision : il a donc servi
à l'entraînement, au sens large. Parmi cinquante essais, le meilleur score
est en partie dû à la chance, et il sera plus flatteur que la réalité. La
bonne pratique consiste à garder un troisième jeu de données, mis de côté
dès le départ, que l'on ne regarde qu'une seule et unique fois, tout à la
fin, pour annoncer le résultat.
\end{solution}
